Model-based agents trade sample efficiency against the risk that the learned model is wrong. We study this trade-off in the smallest setting where it can be measured: tabular Dyna-Q with a last-outcome model on 6x9 mazes whose walls change mid-training (blocking, shortcut) or whose transitions are stochastic (slip \rhval{const/slip_prob:1}), comparing reimplemented Q-learning, Dyna-Q, Dyna-Q+ and prioritized sweeping over 1--100 planning steps per real step, with 20 seeds for the main comparison and 10 for the sweeps. Planning clearly helps when the model is right (Dyna-Q beats Q-learning on the static maze, paired $p=\rhval{cmp/main/q-learning/static/cum_reward/paired_p:sci1}$), and Dyna-Q+ recovers from a stale model and finds a shortcut where plain Dyna-Q does not, with its advantage growing with the number of planning steps. Under slip \rhval{const/slip_prob:1} the benefit of planning disappears: Dyna-Q's mean cumulative reward falls monotonically with $n$, and at $n\ge50$ it is no longer distinguishable from Q-learning (means \rhval{sweep_n/dyna-q@n=50/stochastic/cum_reward/mean:1} and \rhval{sweep_n/dyna-q@n=100/stochastic/cum_reward/mean:1} versus \rhval{sweep_n/q-learning/stochastic/cum_reward/mean:1}, paired $p=\rhval{cmp/h4_n50/dyna-q/stochastic/cum_reward/paired_p:2}$ and $\rhval{cmp/h4_n100/dyna-q/stochastic/cum_reward/paired_p:2}$, 10 seeds). We did not observe the predicted decline of Dyna-Q's post-change reward with $n$ on the blocking maze (the curve is noisy and low at small $n$), and the Dyna-Q+ bonus is fragile: bonus weights $\kappa\ge 0.01$ collapse performance, and the bonus without the untried-action initialisation hurts on the static and stochastic mazes (paired $p\le\rhval{cmp/abl_dynaq_plus/dyna-q+-w-o-untried/stochastic/post_reward/paired_p:sci1}$ against Dyna-Q+ and $p\le\rhval{cmp/abl_ref_dynaq/dyna-q+-w-o-untried/stochastic/post_reward/paired_p:sci1}$ against Dyna-Q), with only non-significant trends elsewhere. Prioritized sweeping was statistically indistinguishable from Dyna-Q in cumulative and post-change reward at $n=10$, though it earned more in the first \rhval{const/reward_window_steps} steps (paired $p=\rhval{cmp/main/prioritized-sweeping/static/early_reward/paired_p:3}$). All claims concern a small CPU study with fixed, untuned hyperparameters.
